\subsection{Artinian Modules and Noetherian Modules}

DEFINITION 1. --- Let A be a ring. We call an A-module M Artinian (resp. Noetherian) if it satisfies the following equivalent conditions:

\begin{enumerate}
\def\labelenumi{(\roman{enumi})}
\item
  Every nonempty set of submodules of M, ordered by inclusion, has a minimal (resp. maximal) element.
\item
  Every decreasing (resp. increasing) sequence of submodules of M is stationary.
\end{enumerate}

The equivalence of conditions (i) and (ii) follows from Set Theory, III, §6, No.~5, p.~190, Proposition 6.

An A-module M is Artinian (resp. Noetherian) if and only if M, viewed as a module over the ring of homotheties \(A_{M}\) , is Artinian (resp. Noetherian).

Let M be an Artinian (resp. Noetherian) A-module. Every nonempty set of submodules of M, ordered by inclusion, that is left directed (resp. right directed) has a least element (resp. a greatest element) (Set Theory, III, §1, No.~10, p.~145, Proposition 10).

Let M be an Artinian (resp. Noetherian) A-module and \((\mathrm{M}_{i})_{i\in\mathrm{I}}\) a family of submodules of M. The intersections (resp. sums) of finite subfamilies of the family \((\mathrm{M}_{i})_{i\in\mathrm{I}}\) form a nonempty left (resp. right) directed set of submodules of M. Therefore, there exists a finite subset J of I such that \(\bigcap_{i\in I}M_{i}=\bigcap_{i\in J}M_{i}\) (resp. \(\sum_{i\in I}M_{i}=\sum_{i\in J}M_{i}\) ).

Examples. --- 1) A finite-dimensional vector space over a field is Artinian and Noetherian.

\begin{enumerate}
\def\labelenumi{\arabic{enumi})}
\setcounter{enumi}{1}
\tightlist
\item
  Let M be an A-module. If there exists an infinite family \((\mathrm{M}_{i})_{i\in\mathrm{I}}\) of nonzero submodules of M whose sum is direct, then M is neither Artinian nor Noetherian: indeed, for every strictly decreasing (resp. strictly increasing) infinite sequence \((\mathrm{J}_{n})\) of subsets of I, the infinite sequence \((\sum_{i\in\mathrm{J}_{n}}\mathrm{M}_{i})\) of submodules of M is strictly decreasing (resp. strictly increasing). In particular, an infinite-dimensional vector space over a field is neither Artinian nor Noetherian.
\end{enumerate}

*3) We will see further on that the Z-module Z is Noetherian but not Artinian (VIII, p.~5, Example 3).*

\begin{enumerate}
\def\labelenumi{\arabic{enumi})}
\setcounter{enumi}{3}
\tightlist
\item
  Let \(p\) be a prime number and \(\mathrm{M}_p\) the \(p\) -primary component of the torsion \(\mathbf{Z}\) -module \(\mathbf{Q} / \mathbf{Z}\) (VII, §2, No.~2, p.~7). Every submodule of \(\mathrm{M}_p\) is equal to either \(\mathrm{M}_p\) or \(p^{-n}\mathbf{Z} / \mathbf{Z}\) for an integer \(n \in \mathbf{N}\) (VII, §2, p.~54, Exercise 3). Consequently, \(\mathrm{M}_p\) is an Artinian but not Noetherian \(\mathbf{Z}\) -module.
\end{enumerate}

PROPOSITION 1. --- An A-module M has finite length (II, §1, No.~10, p.~212) if and only if it is both Artinian and Noetherian.

Suppose that M has finite length d.~Then every strictly increasing or strictly decreasing sequence of submodules of M has at most \(d+1\) terms (I, §4, No.~7, p.~44). Consequently, M is Artinian and Noetherian.

Conversely, suppose that M is Artinian and Noetherian. Let S be the set of submodules of M of finite length. The zero submodule is an element of S, and since M is Noetherian, S has a maximal element N. Let us give a proof by contradiction and suppose that \(M \neq N\) . The set of submodules of M distinct from N and containing N then has a minimal element P because M is Artinian. The module P/N has length 1, and since N is a module of finite length, the same holds for P (II, §1, No.~10, p.~212, Proposition 16). This contradicts the definition of N.

PROPOSITION 2. --- An A-module M is Noetherian if and only if every submodule of M is finitely generated.

Begin by assuming that every submodule of M is finitely generated. Let \((\mathrm{P}_{n})_{n\in\mathbf{N}}\) be an increasing sequence of submodules of M, and let P be its union. This is a submodule of M. By assumption, there exists a finite subset F of M generating the module P; let \(n\in N\) be an integer such that \(F\subset P_{n}\) . We then have \(P_{n}=P\) , and the sequence \((\mathrm{P}_{n})_{n\in\mathbf{N}}\) is stationary. This proves that the module M is Noetherian.

The converse is a consequence of the following more precise statement.

Lemma 1. --- Let M be a Noetherian A-module and E a subset of M. There exists a finite subset F of E generating the same submodule as E.

Indeed, by VIII, p.~2, there exists a finite subset F of E such that \(\sum_{x\in E}Ax=\sum_{x\in F}Ax\) .

PROPOSITION 3. --- Let M be an A-module and N a submodule of M. Then M is Artinian (resp. Noetherian) if and only if N and M/N are.

We will give the proof in the case of Artinian modules; the case of Noetherian modules is analogous.

Suppose that M is Artinian. Since every submodule of N is a submodule of M, the module N is Artinian. Let \((\mathrm{P}_{n})_{n\in\mathbf{N}}\) be a decreasing sequence of submodules of M/N. There exists a decreasing sequence \((\mathrm{Q}_{n})_{n\in\mathbf{N}}\) of submodules of M containing N such that \(P_{n}=Q_{n}/N\) for every \(n\in N\) (I, §4, No.~6, p.~41, Theorem 4). Since M is Artinian, the sequence \((\mathrm{Q}_{n})\) is stationary, hence so is the sequence \((\mathrm{P}_{n})\) . Consequently, the module M/N is Artinian.

Conversely, suppose that the modules N and M/N are Artinian, and consider a decreasing sequence \((\mathrm{P}_{n})\) of submodules of M. The sequence \(P_{n}^{\prime} = N \cap P_{n}\) of submodules of N is stationary. Likewise, the sequence \(\mathrm{P}_{n}^{\prime\prime} = (\mathrm{N} + \mathrm{P}_{n})/\mathrm{N}\) of submodules of M/N is stationary. Hence, there exists an integer \(m \in N\) such that we have \(P_{n}^{\prime} = P_{m}^{\prime}\) and \(P_{n}^{\prime\prime} = P_{m}^{\prime\prime}\) for every integer \(n \geqslant m\) . The sequence \((\mathrm{P}_{n})\) is then stationary by the following lemma.

Lemma 2. --- Let M be an A-module and N, P, and Q submodules of M. Suppose that we have \(P \subset Q\) , \(N \cap P = N \cap Q\) , and \(N + P = N + Q\) . We then have P = Q.

Let x be an element of Q. It belongs to \(N + P\) ; hence, there exists an element y of P such that \(x - y \in N\) . Since Q contains P, the difference x - y belongs to \(N \cap Q\) and therefore to P. Consequently, x belongs to P.

COROLLARY. --- Let M be an A-module and \((\mathrm{M}_{i})_{i\in\mathrm{I}}\) a finite family of submodules of M.

\begin{enumerate}
\def\labelenumi{\alph{enumi})}
\item
  If the modules \(\mathrm{M}_i\) are Artinian (resp. Noetherian), then so is their sum \(\sum_{i\in \mathrm{I}}\mathrm{M}_i\) .
\item
  If the modules \(\mathrm{M} / \mathrm{M}_i\) are Artinian (resp. Noetherian), then so is the module \(\mathrm{M} / \bigcap_{i\in \mathrm{I}}\mathrm{M}_i\) .
\end{enumerate}

By induction, it suffices to treat the case when \(I = \{1, 2\}\) . The module \(M_{2}/(M_{1} \cap M_{2})\) , quotient of \(M_{2}\) , is isomorphic to the submodule \((M_{1} + M_{2})/M_{1}\) of \(M/M_{1}\) (I, §4, No.~6, p.~41, Theorem 4).

In part a), we assume that \(\mathrm{M}_1\) and \((\mathrm{M}_1 + \mathrm{M}_2) / \mathrm{M}_1\) are Artinian (resp. Noetherian); the same then holds for \(\mathrm{M}_1 + \mathrm{M}_2\) (Proposition 3).

In part b), we assume that \(M/M_{2}\) and \(M_{2}/(M_{1} \cap M_{2})\) are Artinian (resp. Noetherian); the same then holds for \(M/(M_{1} \cap M_{2})\) (loc. cit.).

Example 5. --- Let \((\mathrm{M}_{i})_{i\in\mathrm{I}}\) be a finite family of A-modules. If the modules \(M_{i}\) are Artinian (resp. Noetherian), then so is their direct sum \(\bigoplus_{i\in I} M_{i}\) .

Remark. --- The definitions and results of this subsection extend to arbitrary abelian groups with operators by replacing the submodules in the statements with stable subgroups.

